\chapter{ΣΥΜΠΕΡΑΣΜΑΤΑ}

\section{Ανασκόπηση της Μελέτης}

Στην συγκεκριμένη διπλωματική εργασία αναπτύχθηκε, εφαρμόστηκε και αξιολογήθηκε μια ολοκληρωμένη μεθοδολογία για την εκτίμηση, ταυτοποίηση και συσταδοποίηση των ηλεκτρομηχανικών ταλαντώσεων σε Συστήματα Ηλεκτρικής Ενέργειας (ΣΗΕ). Γίνεται μια προσπάθεια αντιμετώπισης των σύγχρονων προκλήσεων, που συνεπάγεται με τη μείωση της φυσικής αδράνειας του δικτύου λόγω της αυξημένης διείσδυσης Ανανεώσιμων Πηγών Ενέργειας (ΑΠΕ). Στο προαναφερθέν πλαίσιο, η εφαρμογή της μεθόδου Matrix Pencil (MP) σε σήματα Ringdown μετά από σοβαρές διαταραχές επιβεβαίωσε την υψηλή πιστότητα ανακατασκευής των κυματομορφών τάσης, ρεύματος, ενεργού και άεργου ισχύος, επιδεικνύοντας εξαιρετική ακρίβεια στην εξαγωγή των συχνότητων ταλάντωσης και των συντελεστών απόσβεσης, όπως επαληθεύτηκε από τους δείκτες $R^2$, RMSE και MAD. Παράλληλα, στην Ambient ανάλυση, η τεχνική αναγνώρισης υποχώρου N4SID αποδείχθηκε ιδιαίτερα αποτελεσματική στην εκτίμηση των δεσποζόντων τρόπων ταλάντωσης υπό συνθήκες φυσιολογικής λειτουργίας (ambient θόρυβος), καταφέρνοντας να εντοπίσει τους κρίσιμους ενδο-περιοχικούς και δια-περιοχικούς ρυθμούς ταλάντωσης παρά την παρουσία θορύβου μετρήσεων.

Όσον αφορά τη συγκριτική αξιολόγηση των αλγορίθμων συσταδοποίησης, η διερεύνηση επτά διαφορετικών τεχνικών (k-Means, k-Medoids, DBSCAN, OPTICS, HDBSCAN, GMM, Agglomerative Hierarchical Clustering) έδειξε ότι οι αλγόριθμοι που βασίζονται στην πυκνότητα (DBSCAN, HDBSCAN, OPTICS) καθώς και τα μίγματα Γκαουσιανών κατανομών (GMM) υπερέχουν στη διαχείριση θορύβου και στην απόρριψη αριθμητικών ή ψευδών ιδιοτιμών. Η υψηλή συμφωνία των παραγόμενων συστάδων με τις τιμές αναφοράς επιβεβαιώθηκε και ποσοτικά μέσω των εξωτερικών δεικτών ARI και V-measure.

Η μελέτη της επίδρασης των αιολικών μονάδων έδειξε ότι η αντικατάσταση συμβατικών σύγχρονων γεννητριών από αιολικά πάρκα τεχνολογίας DFIG στο σύστημα IEEE 39-Bus προκάλεσε ορατή μετατόπιση των ηλεκτρομηχανικών ρυθμών, επηρεάζοντας τόσο τη συχνότητα όσο και την απόσβεση των δια-περιοχικών ταλαντώσεων. Το προτεινόμενο πλαίσιο ταυτοποίησης και συσταδοποίησης προσαρμόστηκε με επιτυχία στις νέες συνθήκες, αποδεικνύοντας την αξιοπιστία και τη στιβαρότητά του σε περιβάλλον ηλεκτρικών δικτύων με χαμηλή αδράνεια.

\section{Συμβολή της Διπλωματικής}

Η κύρια συμβολή της διπλωματικής εργασίας έγκειται στην ανάπτυξη ενός ολοκληρωμένου και συστηματικού πλαισίου ταυτοποίησης και συσταδοποίησης για τη διαχείριση δεδομένων WAMS. Η προτεινόμενη μεθοδολογία συνδυάζει προηγμένες τεχνικές ταυτοποίησης στο πεδίο του χρόνου και του υποχώρου με επτά διαφορετικούς αλγόριθμους συσταδοποίησης, προσφέροντας μια ολιστική προσέγγιση στην ανάλυση των ηλεκτρομηχανικών ταλαντώσεων.

Παρέχεται, επίσης, μια εκτενής συγκριτική αξιολόγηση αλγορίθμων συσταδοποίησης (k-Means, k-Medoids, DBSCAN, OPTICS, HDBSCAN, GMM, Agglomerative Hierarchical Clustering) ως προς τη συνοχή, την κάλυψη, τη στιβαρότητα απέναντι στον θόρυβο και την ακρίβεια χαρτογράφησης των ιδιοτιμών. Για την αντικειμενική μέτρηση της απόδοσης των αλγορίθμων αυτών σε σχέση με τις θεωρητικές τιμές αναφοράς, εισάγεται η συστηματική χρήση εξειδικευμένων δεικτών αξιολόγησης, όπως ο Median Absolute Deviation (MAD), ο Adjusted Rand Index (ARI) και το V-measure.

Η εργασία συμβάλλει στη μελέτη της επίπτωσης της αυξημένης διείσδυσης ΑΠΕ στη δυναμική του δικτύου, αποτυπώνοντας ποσοτικά τη μεταβολή των ηλεκτρομηχανικών ρυθμών λόγω της αντικατάστασης συμβατικών γεννητριών από αιολικά πάρκα τεχνολογίας DFIG. Παράλληλα, αξιολογείται η προσαρμοστικότητα και η αξιοπιστία των μεθόδων συσταδοποίησης υπό τις μεταβαλλόμενες συνθήκες λειτουργίας του συστήματος.

\section{Προτάσεις για περαιτέρω Έρευνα}

Με βάση τη μεθοδολογία και τα ευρήματα της παρούσας διπλωματικής εργασίας, προτείνονται οι ακόλουθες κατευθύνσεις για μελλοντική έρευνα:

\begin{enumerate}
    \item Εφαρμογή σε Πραγματικό Χρόνο: Επέκταση των αλγορίθμων συσταδοποίησης για τη συνεχή επεξεργασία ροών δεδομένων PMU σε πραγματικό χρόνο, με ενσωμάτωση μηχανισμών αυτόματης προσαρμογής παραμέτρων υπό δυναμικά μεταβαλλόμενες συνθήκες δικτύου.
    \item Δοκιμή σε Πραγματικές Μετρήσεις WAMS: Επαλήθευση της μεθοδολογίας με χρήση πραγματικών πεδίων μετρήσεων από ηλεκτρικά δίκτυα μεταφοράς, εξετάζοντας την ανθεκτικότητα των αλγορίθμων σε φαινόμενα απώλειας δεδομένων ή λανθασμένων μετρήσεων.
    \item Ενσωμάτωση Τεχνικών Βαθιάς Μάθησης: Συνδυασμός των παραδοσιακών μεθόδων αναγνώρισης υποχώρου με αρχιτεκτονικές νευρωνικών δικτύων για την ταχύτερη και πιο αποτελεσματική εκτίμηση ιδιοτιμών σε πολύπλοκα δίκτυα ευρείας κλίμακας.
\end{enumerate}
